\title{瓶颈 Adapter 如何影响规则生成？\\激活函数、宽度与插入位置的对照研究}
\PaperAuthor
\maketitle

\begin{abstract}
给定一份中文需求，语言模型可以填对当前记录，却为今后的记录写下错误的验证规则。我们检验瓶颈 Adapter 能否只约束指定字段：在 Qwen3-4B 上对照激活函数、瓶颈宽度与插入位置，以未微调模型和 LoRA 为参照，用固定记录、每次只改一条规则的配对请求计分。将宽度为 16 的仿射模块从解码器块末移到注意力残差之后，整对全部通过的题对由 3/16 增至 8/16，而两组对目标规则的判定同为 30/32；一个案例把差异定位到块末模块给无关数组写下的无根据上限。这个单种子实验说明，判断适配效果需要同时问：该改的规则改对了吗？不该改的规则保留了吗？
\end{abstract}

\section{引言}
设想一个运行在个人设备上的本地小模型助理，接到“帮我定个闹钟”的请求。闹钟应用只接受严格结构化的记录：字段名、类型或外层包装出错即被拒绝。零样本时，这一规模的模型往往恰好卡在这一步。常规补救是冻结一个骨干、按应用训练一个可插拔的小模块，这正是瓶颈 Adapter 面向的场景 \citep{houlsby2019}。

然而记录只是交付物的一半。同一句请求还陈述了\emph{未来}记录的规则：哪些字段必填、备注可否为空、重复列表最多包含几项。我们要求模型把这些规则与当前实例一并写成可执行的 JSON Schema，而之后的每条记录都依据这份契约验证，而非依据今天的例子。Schema 因此是为尚未提交的记录编写的程序：今天的闹钟填对了，并不能保证下个月一条合法的闹钟会被接受。最有破坏性的失败是沉默而延迟的，例如一个没人约束的数组被加上了长度上限，在生成之后很久才开始拒绝合法记录。本文实验使用合成中文需求，其规则族——数量界限、可选性、可空性、条件依赖——与这类助手侧约束同构；我们不在闹钟数据上做评估。

Adapter 设计在三个轴上不同：瓶颈内的激活函数、瓶颈宽度，以及更新施加于 transformer 的位置。我们把这三个轴当作待检验的假设而非沿袭的默认。我们的问题是：当一条需求限制一个数组而放任另一个时，哪些设计能把界限写在目标数组上、并让另一个保持自由？在固定实例、每次只改变一条规则句的配对测试中，插入位置产生了最清晰的分离——移动同一个窄模块，其对变更规则的判定保持不变，整对全部通过的数目显著增加，逐例检查把差异定位到写在无关数组上的一个无根据的长度上限。本文的贡献有三：（i）把变更规则正确性与不变约束保真分开计分的可执行配对评估；（ii）在现代解码器大模型上预指定的激活、宽度与位置单变量对照；（iii）一个逐例追踪的案例，显示相同的规则判定总数可以掩盖截然不同的限制行为。

\section{相关工作}
瓶颈 Adapter 在冻结网络上附加降维再升维的小模块 \citep{houlsby2019}，此后放置位置一直是设计变量：前馈子层后放置可换用、复用的语言与任务模块 \citep{pfeiffer2020}，不重训任何模块即可组合已训练的任务 Adapter \citep{pfeiffer2021}，低层模块可删而代价很小 \citep{ruckle2021}。\citet{he2022} 把这些变体统一到两个可分离的轴上——更新如何计算、施加于何处——并发现与前馈子层并行接入的瓶颈模块是特别有效的放置。在这条工作线上，放置一直按“每可训练参数的任务准确率”挑选；放置对适配后模型所写内容的影响——施加哪些约束、施加在哪——尚未被单独分离。我们的 D--E 对照正是做这件事：在解码器大模型上沿位置轴移动同一训练形态，度量生成的规则行为。

可解释性工作为位置为何重要提供了一个视角。在 transformer 中，各子层通过共享的残差流通信：每个注意力与前馈模块从残差流读取，并向其写回 \citep{elhage2021}。插入位置因此决定了哪些计算能够消费这次更新：放在注意力之后、本块前馈网络之前的模块会改变该网络的输入，若其输出同时被向前携带，后续各层也会读到；放在完整块之后的模块则只被后续层读取。我们将其作为设计假设而非机理断言：本文的度量停留在行为与结构层面，并未分析训练出的模块内部实现了什么。

行为测试围绕特定能力、不变性和方向性变化组织评估 \citep{ribeiro2020}。我们将这一思路用于可执行的 JSON Schema 验证决策。Draft 2020-12 定义了验证语义 \citep{jsonschema}：即使数值正确，附着在错误模式节点上的关键词也未必能约束目标数组。与模式字符串匹配或仅验证给定实例相比，成对要求及合法替代实例能更直接地揭示这类错误。

\section{实验设计}
\paragraph{目标与设计依据。}
本研究只回答一个问题：当一条需求给其中一个列表加上数量上限、另一个列表不加限制时，哪种 Adapter 设计能把上限写在被点名的列表上、而不去动另一个？回答它需要三样东西。其一，每组只差一个设计点的模型：我们在原训练设置下训练三个新 Adapter 变体，事先定好 A--B（激活）、B--D（宽度）、D--E（插入位置）三组对比。其二，直接检查规则本身的测试。总分做不到这一点：把未适配基座本来就输出的两个对象合并成要求的外层对象、内容一字不改，总分就能从 20.9\% 抬到 35.2\%；早期 pilot 也表明，Adapter 与 LoRA 训练能把总分从 20.9\% 提高到 36.0\% 与 48.1\%（580 条从未见过的测试题），却说不清分数涨在哪里。把生成的规则和参考答案逐字比对同样不可靠；只检查题目里那条记录，规则本身则完全没有被考到。其三，两个分开的分数：变更规则正确性——被改动的那条要求有没有写进规则里；以及不该改的规则是否原样保留。这遵循行为测试“每次只改一处”的出题思路 \citep{ribeiro2020}。

\paragraph{假设与条件。}
激活：在本身已非线性的解码器内部，非线性瓶颈可能是冗余的；A--B 检验同宽度下 SiLU 是否改变规则行为。宽度：瓶颈限制了更新的秩；B--D 检验更窄的更新能否在保住目标规则的同时收窄其他影响。位置：残差流视角预言，被本块前馈网络读到的更新与只被后续层读到的更新作用范围可能不同；D--E 检验这一差异是否体现在约束落点上。
全部条件均使用 Qwen/Qwen3-4B-Instruct-2507 \citep{qwenmodel}。表~\ref{tab:models} 列出六个实验组。A 为宽度 64 的 SiLU Adapter；B 移除其激活函数；D 将 B 缩窄至宽度 16；E 将 D 移至注意力后残差。C 在注意力 q/k/v/o 投影上使用秩 16 的 LoRA \citep{hu2022}，参数 $\alpha=32$，有效缩放系数 $\alpha/r=2$，无偏置和 dropout。预先指定的局部对比为 A--B、B--D 和 D--E；A--C 作为配置层面的参照。

\ModelTable

对于 Adapter 输入 $h$，两类变换分别为
\begin{align}
T_{\sigma}(h)&=h+U\sigma(Dh+b_D)+b_U,\\
T_{\mathrm{aff}}(h)&=h+U(Dh+b_D)+b_U.
\end{align}
A 使用 $\sigma=\operatorname{SiLU}$；B/D/E 使用仿射形式。其更新的线性部分秩不超过瓶颈宽度，而外围解码器仍是非线性的。A/B 的初始张量相同，D/E 也相同。下投影权重从标准差为 0.02 的正态分布初始化，下投影偏置与上投影均置零，使每个模块初始时均为恒等变换。B/D 的张量形状不同。

E 同时改变传递的残差和 MLP 输入。设 $r$ 为传入块的残差与注意力输出之和，则
\begin{equation}
\begin{gathered}
z=T_{\mathrm{aff}}(r),\\
h_{\mathrm{out}}=z+\operatorname{MLP}(\operatorname{RMSNorm}(z)).
\end{gathered}
\end{equation}
D 则变换完整解码器块的输出。小模块测试在 FP32/BF16、有无检查点重计算的条件下，核对了 E 的输出、输入梯度及 Adapter 梯度与该方程的一致性。

\paragraph{训练与选择。}
五个微调组采用相同的 2,300 条合成训练请求、339 条验证请求、随机种子 42 和三轮训练策略。训练使用 AdamW，学习率 $10^{-4}$，权重衰减 0.01，余弦衰减，5\% 预热，微批量为 1，梯度累积 16 步，梯度范数裁剪阈值为 1.0，序列上限为 4,096 token。损失仅计算助手目标 token。基座以 BF16 和 SDPA 计算；Adapter 以 FP32 计算，再将输出转回输入数据类型。训练使用 NVIDIA L40S 和 PyTorch 2.7.1+cu128。

每轮训练后解码验证集，按联合成功率、数据准确率、平衡规则得分、自一致性的字典序选择检查点；完全相同时保留较早轮次。各组选中的轮次见表~\ref{tab:models}。贪心推理使用相同的原始提示和聊天模板，不提供示例，批量为 16，生成上限为 2,048 token。

\paragraph{冻结的成对测试。}
旧测试集包含 48 对任务，覆盖计数边界、可选性、可空性和条件依赖，各有两种对齐措辞。确认集包含 24 对新任务，其中 16 对为计数任务，8 对为可空性对照。每对仅改变一个规则句，事实与参考数据保持不变。计数区间相交但互不包含：$[2,4]$ 对 $[1,3]$，或 $[3,5]$ 对 $[2,4]$；变化方向与字段顺序均衡。指定目标数组承载被操纵的边界，第二个数组则明确不限制长度。共享参考数据同时满足两臂要求。

确认集分为原始、改写和重排版本。改写在原位置重写语义块；重排仅置换原始句子的顺序。探针标签、协议和分析计划均在查看预测前冻结。探针覆盖目标数组长度 0 至 6、干扰数组长度 0 和 6，以及必填性、类型、额外字段和材料值。

\paragraph{共同主要指标。}
令 $c_{iap}$ 表示任务对 $i$ 的臂 $a$ 是否正确判定探针 $p$。变更探针正确性对区分两臂的比较取 $c_{i0p}c_{i1p}$ 的平均；整对通过率对任务对取 $\prod_{a,p}c_{iap}$ 的平均。计数任务的分母分别为 32 次比较和 16 对任务。缺失、格式错误及不可执行的回复均保留为失败。条件变更正确性仅纳入两臂均可执行的比较，并报告实际分母。评分时严格检查外层对象，并按 Draft 2020-12 执行不含引用的规则。可以省略方言声明；显式声明其他方言则计为失败。整对通过衡量的是规则行为；严格联合交付还要求数据正确、符合接口。对两个裸对象进行对称包装，不改变本机制研究的套件计数。全部预测均已到齐，确认集无截断回复。

\section{结果}
\ResultsTable

\subsection{激活函数与宽度改变错误分布}
如表~\ref{tab:results} 所示，移除 SiLU（A--B）后，原始版本的变更探针正确数仍为 29/32，改写版本从 31/32 变为 30/32；重排版本则从 19/32 提高至 30/32。A 在五个重排样例中仅将准确边界放在对象根节点，B 没有此类样例。然而，原始/改写/重排版本的数量题整对通过数从 9/13/6 降至 8/8/5。修正一种放置错误，并不确保整体模式行为更好。

缩窄仿射模块（B--D）后，数量题整对通过数从 8/8/5 变为 3/8/3。两者在每个版本的全部 32 个样例中，都将准确边界直接放在目标数组上。因此，这一差异不能概括为目标区间丢失：D 在 22/14/23 个样例中限制了无限制的干扰数组，B 则为 14/14/22。可空性套件成功数反而从 5/2/7 提高至 8/5/8，各版本分母均为 8。可见，缩小瓶颈并非使每类规则都变差。

\subsection{插入位置改善不变约束，变更规则正确数相同}
将 D 移至注意力后残差（E），数量题整对通过数从 3/8/3 提高至 8/9/6。E 保留了 D 的每一个成功任务对，并分别新增 5、1、3 对。两者的变更探针正确数完全相同，均为 30/29/28，分母均为 32。条件分母区分了重排版本的结果：在可执行比较中，D 为 28/30，E 为 28/32。无条件计数将这些失败保留在内。

错误地给干扰数组添加长度限制的输出数从 22/14/23 降至 14/11/16，而目标边界的准确放置始终保持为 32/32。D/E 在原始和重排版本中也均通过全部 8 对可空性套件，改写版本分别通过 5 对和 6 对。因此，这一局部位置变更最清楚的收益，是更完整地保留约束，而非提高对变更区间的判定正确数。

在原句版本的题对 \texttt{confirm\_count\_04} 中，两模型都正确地将 \texttt{beam\_presets} 在 a 臂限制为 $[2,4]$，在 b 臂限制为 $[3,5]$。除 D 在两臂的 \texttt{gel\_codes} 上添加了 \texttt{maxItems: 2} 外，其余模式相同。该数组明确不限制长度。未变更的六元素干扰探针在两臂下都合法：D 拒绝它，E 接受它。这是 D 每臂唯一失败的探针，而 E 通过两臂的完整套件。该案例直接定位了变更规则得分相同背后的非预期限制。

\subsection{旧任务的成功不能原样迁移}
在 Base 两种措辞下均通过的 10 对旧计数任务中，A 没有一对能在两个版本上同时保留成功，B 保留 3 对，D 保留全部 10 对，E 保留 9 对。D 在旧任务的每种措辞下均通过全部 12 对计数任务，却仅有 2/16 对确认任务能在三个版本下全部通过。新增干扰数组和新场景揭示了旧任务排序掩盖的失败。

LoRA C 在确认集通过 11/13/8 对计数套件，其中 8 对在三个版本中均成功；A/B/D/E 的跨版本成功数则为 6/4/2/5。C 同时改变了更新目标、参数化、缩放和偏置。Base 自身的波动也很大：各版本通过 2/12/1 对套件，仅有一对在全部版本中通过。没有任何适配模型在全部版本中保留这一对。

\section{结论}
Adapter 架构既改变约束落在何处，也改变会出现哪些无关限制。注意力后的窄仿射 Adapter 在变更规则正确数与块后版本相同的情况下，提高了整对通过率。移除激活函数和缩窄宽度则产生各不相同、优劣混合的错误分布。多种子重复与其他规则族将检验这一位置效应能推广多远。对生成的验证器，应将指定规则变更、未变更的合法替代实例以及对齐的措辞/句序版本共同评估：区分性探针上的高分，可以与模式其他位置上的有害限制同时存在。

\section*{局限性}
瓶颈 Adapter 方法本身存在边界。每条轴的收益都伴随代价：去掉 SiLU 修复了根节点放置错误，却没有提高整对通过数；收窄宽度保住了目标界限，却收紧了对无关数组的限制；位置迁移收益最大，但最好的配置在三个版本中仍分别留下 8、7、10 对计数题未通过（分母各为 16）。在相近的参数预算下，LoRA 在早期实验的联合成功率上超过当时的 Adapter，在本研究的确认集整对通过数上也追平或超过所有 Adapter 变体，说明这一 Adapter 家族在本任务上的上限低于已有的替代方案。原因或可一言蔽之：附挂于冻结骨干的低秩残差更新只能调制既有表征，与骨干自身任务理解绑定的规则错误不在其可达范围之内。此外，有限探针确定的是受测决策而非模式等价性，因此可执行得分与结构放置计数并列报告。实验方面的不足：每个配置仅一个种子和一个经验证集选择的轮次；确认集规模较小（24 对），各措辞版本是对齐观测而非独立样本；早期实验的冻结得分在参考标签与要求不一致处衡量的是与参考见证的相符程度。

\section*{AI 与数据披露}
全部数据均为合成数据。原始生成器借助 GPT-6 Luna 和确定性替换编制。不存在独立的人工标注测试集，但数据经过人工抽样检查以验证正确性。参考模式和探针通过程序验证。结果来自保留的预测和确定性分析；冻结协议、模型版本、提示、生成设置和逐项结果均已上传至 GitHub 留档。\footnote{\url{https://github.com/Jianqu0817/NLP_FINAL_DATA}}

